\section{Conclusion}\label{sec:conclusion}

Competition for mobile production can generate excessive local protection when public reliability and privately chosen contingency readiness interact. In a nonempty open set of primitives within the specified model, the coordinated protection-policy planner chooses zero additional protection, whereas decentralized governments choose strictly positive protection in a symmetric Nash equilibrium. The conflict arises because improved primary-site reliability can crowd out privately financed backup readiness sufficiently to reduce national industrial surplus, even while the resulting attraction of primary plants benefits an individual jurisdiction.

The analysis also identifies a distinct comparative-static implication. When protection generates an additional direct benefit valued by governments and the planner, decentralized overprovision can persist with a strictly positive coordinated optimum. In an exactly solved separable-product benchmark, increasing this benefit widens the decentralized--coordinated protection gap when both readiness and primary location are endogenous, leaves the gap unchanged when readiness is fixed, and narrows it when locations are fixed. Each restriction is imposed throughout its policy game, including unilateral deviations. The positive full-response gap effect persists locally with strictly positive product substitutability.

These results concern coordination of public protection while firms continue to choose their locations and resilience privately. They do not establish a universal ranking across technologies or a comprehensive optimum for climate adaptation, because the baseline welfare measure excludes protection benefits outside the modeled industry. The central implication for public-investment appraisal is to account jointly for direct reliability improvements, induced changes in private contingency readiness, and competition for mobile production.
